\documentclass{article}
\usepackage{hyperref}
\usepackage{xcolor}
\usepackage{graphicx}

\title{\Huge \textbf{Telegram Bot \& Private Backend for Appointment Monitoring}}
\author{}
\date{}

\begin{document}

\maketitle

\section*{\textcolor{red}{⚠ Disclaimer}}
This project is created \textbf{for educational purposes only}. \textbf{Commercial use is strictly prohibited} and may be \textbf{illegal} depending on your jurisdiction. The author is \textbf{not responsible} for any misuse of this software.

\section*{📌 Features}
\begin{itemize}
    \item ✅ Telegram bot updates a message in a channel periodically.
    \item ✅ Backend fetches appointment data securely.
    \item ✅ Bot deletes old messages and sends new updates only when needed.
    \item ✅ Dockerized for easy deployment.
    \item ✅ Backend is private and only accessible inside Docker.
\end{itemize}

\section*{📽 Demo Video}
\begin{center}
    \href{https://your-video-link-here}{\includegraphics[width=0.8\textwidth]{https://via.placeholder.com/800x450?text=Demo+Video}}
\end{center}

\section*{🚀 Installation \& Usage}

\subsection*{1️⃣ Clone the Repository}
\begin{verbatim}
git clone https://github.com/yourusername/telegram-appointment-bot.git
cd telegram-appointment-bot
\end{verbatim}

\subsection*{2️⃣ Configure \texttt{.env} Files}
Create \texttt{.env} files inside both \textbf{backend} and \textbf{bot} directories.

\subsubsection*{Backend (\texttt{/backend/.env})}
\begin{verbatim}
PORT=3000
EXAM_TYPE=G
LAST_NAME=YourLastName
LICENSE_NUMBER=123456
\end{verbatim}

\subsubsection*{Bot (\texttt{/bot/.env})}
\begin{verbatim}
TELEGRAM_BOT_TOKEN=your_bot_token
TELEGRAM_CHANNEL_ID=your_channel_id
API_URL=http://backend:3000/appointments  # Use service name inside Docker
\end{verbatim}

\subsection*{3️⃣ Build \& Run Using Docker}
\begin{verbatim}
docker-compose up -d --build
\end{verbatim}

\textbf{Note:} The bot will \textbf{start fetching and updating appointments}, and the backend will run \textbf{privately inside Docker}.

\section*{📜 API Endpoints (Backend)}
\begin{center}
\begin{tabular}{|c|l|l|}
\hline
\textbf{Method} & \textbf{Endpoint} & \textbf{Description} \\
\hline
GET & \texttt{/appointments} & Fetch all available appointments \\
GET & \texttt{/appointments/:locationId} & Fetch appointments for a specific location \\
GET & \texttt{/locations} & Get available locations \\
\hline
\end{tabular}
\end{center}

\section*{🛠 Development}

\subsection*{Backend}
\begin{verbatim}
cd backend
npm install
npm run dev
\end{verbatim}

\subsection*{Bot}
\begin{verbatim}
cd bot
npm install
npm run dev
\end{verbatim}

\section*{📜 Stopping \& Removing the Containers}
\begin{verbatim}
docker-compose down
\end{verbatim}

\section*{⚠ Legal Disclaimer}
This project is \textbf{strictly for educational purposes}. \\
Using this bot for \textbf{commercial purposes} is \textbf{illegal} and may \textbf{violate terms of service}. \\
The author \textbf{assumes no responsibility} for any misuse of this software.

\section*{📜 License}
This project is \textbf{open-source} but \textbf{not for commercial use}.

\section*{🤝 Contributing}
Pull requests are welcome! Feel free to \textbf{fork this repo} and suggest improvements. 🚀

\end{document}